% ─────────────────────────────────────────────────────────────────────
%   ANNEXES
% ─────────────────────────────────────────────────────────────────────

% Configuration globale des listings de code
\lstset{
  basicstyle=\ttfamily\footnotesize,
  breaklines=true,
  breakatwhitespace=false,
  columns=fullflexible,
  keepspaces=true,
  showstringspaces=false,
  tabsize=2,
  frame=single,
  rulecolor=\color{ensamBlue!40},
  backgroundcolor=\color{ensamBlue!3},
  xleftmargin=0.5cm,
  xrightmargin=0.5cm,
  keywordstyle=\color{ensamBlue}\bfseries,
  stringstyle=\color{orange!70!black},
  commentstyle=\color{gray}\itshape,
  numberstyle=\tiny\color{gray},
  numbers=left,
  numbersep=6pt,
  captionpos=b,
}

\chapter{Annexes}

% =====================================================================
\section{Extraits de code source}

Cette annexe présente les extraits de code source les plus représentatifs de
l'application développée pour CIVCO BTP, sélectionnés pour illustrer les
contributions techniques décrites dans ce rapport.

% ─────────────────────────────────────────────────────────────────────
\subsection{Authentification API (\texttt{AuthController.php})}

Le contrôleur d'authentification gère la connexion des utilisateurs internes
via Laravel Sanctum. La méthode \texttt{login} valide les identifiants, régénère
la session et retourne le contexte utilisateur (profil, rôles, permissions).

\begin{lstlisting}[language=PHP,
  caption={AuthController.php : connexion utilisateur (extrait)}]
public function login(LoginRequest $request,
    AuthContextService $authContext): JsonResponse
{
    $credentials = $request->only('email', 'password');

    if (! Auth::attempt($credentials, $request->boolean('remember'))) {
        throw ValidationException::withMessages([
            'email' => [__('auth.failed')],
        ]);
    }

    $user = Auth::user();

    if (! $user->canAccessApplication()) {
        Auth::logout();
        throw ValidationException::withMessages([
            'email' => [CheckUserStatus::DEACTIVATED_MESSAGE],
        ]);
    }

    if ($request->hasSession()) {
        $request->session()->regenerate();
    }

    return response()->json($authContext->forUser($user));
}
\end{lstlisting}

% ─────────────────────────────────────────────────────────────────────
\subsection{Conversion devis vers facture (\texttt{QuoteService.php})}

Le service métier \texttt{QuoteService} encapsule la logique de conversion d'un
devis accepté en facture. L'opération s'exécute dans une transaction base de
données pour garantir la cohérence des enregistrements.

\begin{lstlisting}[language=PHP,
  caption={QuoteService.php : conversion devis $\rightarrow$ facture (extrait)}]
public function convertToInvoice(Quote $quote,
    ?int $dispatchNoteId = null): Invoice
{
    if ($quote->status !== QuoteStatus::Accepted) {
        throw new InvalidArgumentException(
            'Only accepted quotes can be converted.');
    }

    if ($quote->invoice()->exists()) {
        throw new InvalidArgumentException(
            'This quote has already been converted.');
    }

    return DB::transaction(function () use ($quote, $dispatchNoteId) {
        $quote->load('lines');

        $invoice = Invoice::query()->create([
            'company_id' => $quote->company_id,
            'client_id'  => $quote->client_id,
            'project_id' => $quote->project_id,
            'quote_id'   => $quote->id,
            'reference'  => $this->invoiceReferenceService
                                ->nextForCompany($quote->company_id),
            'status'     => InvoiceStatus::Draft,
            'total_ht'   => $quote->total_ht,
            'total_tax'  => $quote->total_tax,
            'total_ttc'  => $quote->total_ttc,
        ]);

        foreach ($quote->lines as $index => $line) {
            $invoice->lines()->create([
                'sort_order'      => $index + 1,
                'description'     => $line->description,
                'quantity'        => $line->quantity,
                'unit_price_ht'   => $line->unit_price_ht,
                'line_total_ht'   => $line->line_total_ht,
                'line_total_ttc'  => $line->line_total_ttc,
            ]);
        }

        return $invoice->load(['client', 'project', 'lines']);
    });
}
\end{lstlisting}

% ─────────────────────────────────────────────────────────────────────
\subsection{Modèle Projet (\texttt{Project.php})}

Le modèle Eloquent \texttt{Project} centralise les relations métier d'un chantier :
client, phases, équipe, devis et factures.

\begin{lstlisting}[language=PHP,
  caption={Project.php : relations Eloquent (extrait)}]
class Project extends Model
{
    protected $fillable = [
        'company_id', 'client_id', 'reference', 'title',
        'status', 'budget', 'progress_percent',
        'start_date', 'end_date', 'latitude', 'longitude',
    ];

    protected function casts(): array
    {
        return [
            'status' => ProjectStatus::class,
            'start_date' => 'date',
            'end_date' => 'date',
            'budget' => 'decimal:2',
            'progress_percent' => 'decimal:2',
        ];
    }

    public function client(): BelongsTo
    {
        return $this->belongsTo(Client::class);
    }

    public function phases(): HasMany
    {
        return $this->hasMany(ProjectPhase::class);
    }

    public function teamMembers(): BelongsToMany
    {
        return $this->belongsToMany(User::class, 'project_user');
    }

    public function quotes(): HasMany
    {
        return $this->hasMany(Quote::class);
    }
}
\end{lstlisting}

% ─────────────────────────────────────────────────────────────────────
\subsection{Client HTTP frontend (\texttt{client.js})}

Le module \texttt{api/client.js} configure l'instance Axios partagée par tous
les modules frontend. Il transmet les cookies Sanctum, le jeton CSRF et le
contexte société (\texttt{X-Company-Id}).

\begin{lstlisting}[language=Java,
  caption={client.js : configuration Axios et intercepteurs (extrait)}]
const api = axios.create({
  baseURL: import.meta.env.VITE_API_URL ?? '',
  withCredentials: true,
  headers: {
    Accept: 'application/json',
    'Content-Type': 'application/json',
    'X-Requested-With': 'XMLHttpRequest',
  },
})

api.interceptors.request.use(async (config) => {
  if (activeCompanyId) {
    config.headers['X-Company-Id'] = activeCompanyId
  }

  const method = config.method?.toLowerCase()
  if (['post', 'put', 'patch', 'delete'].includes(method)) {
    let token = getCookie('XSRF-TOKEN')
    if (!token) {
      await api.get('/sanctum/csrf-cookie')
      token = getCookie('XSRF-TOKEN')
    }
    if (token) {
      config.headers['X-XSRF-TOKEN'] = token
    }
  }

  return config
})

api.interceptors.response.use(
  (response) => response,
  (error) => {
    if (error.response?.status === 401) {
      window.dispatchEvent(new CustomEvent('auth:unauthorized'))
    }
    return Promise.reject(error)
  }
)
\end{lstlisting}

% =====================================================================
\section{Configuration : variables d'environnement}

Le fichier \texttt{.env} du backend centralise la configuration de l'application.
Le Tableau~\ref{tab:annexe_env} décrit les variables principales utilisées en
développement et en production.

\begin{table}[H]
\caption{Variables d'environnement principales}
\label{tab:annexe_env}
\centering
\small
\begin{tabular}{|l|p{8.5cm}|}
\hline
\textbf{Variable} & \textbf{Rôle} \\
\hline
\texttt{APP\_URL} & URL publique du backend Laravel \\
\hline
\texttt{DB\_CONNECTION} & Pilote SGBD (\texttt{mysql} ou \texttt{pgsql}) \\
\hline
\texttt{DB\_DATABASE} & Nom de la base de données \\
\hline
\texttt{SESSION\_DRIVER} & Stockage des sessions (\texttt{database} recommandé) \\
\hline
\texttt{SANCTUM\_STATEFUL\_DOMAINS} & Domaines autorisés pour les cookies SPA \\
\hline
\texttt{FRONTEND\_URL} & URL du frontend React (CORS et redirections) \\
\hline
\texttt{MAIL\_*} & Configuration SMTP pour les notifications \\
\hline
\end{tabular}
\end{table}

\medskip

Exemple de configuration minimale pour un environnement de développement local :

\lstset{
  language={},
  keywordstyle=\color{black},
  stringstyle=\color{black},
  commentstyle=\color{gray}\itshape,
}

\begin{lstlisting}[caption={.env : configuration de développement (extrait)}]
APP_NAME="CIVCO BTP"
APP_ENV=local
APP_DEBUG=true
APP_URL=http://127.0.0.1:8000

DB_CONNECTION=mysql
DB_HOST=127.0.0.1
DB_PORT=3306
DB_DATABASE=civco_btp
DB_USERNAME=root
DB_PASSWORD=

SESSION_DRIVER=database
SESSION_LIFETIME=120

SANCTUM_STATEFUL_DOMAINS=localhost,127.0.0.1:5173

FRONTEND_URL=http://127.0.0.1:5173

MAIL_MAILER=smtp
MAIL_HOST=smtp.example.com
MAIL_PORT=587
MAIL_USERNAME=null
MAIL_PASSWORD=null
\end{lstlisting}

% =====================================================================
\section{Comptes de démonstration}

Le seeder de démonstration fournit des comptes utilisateurs permettant de
tester les différents profils métier de CIVCO BTP. Le Tableau~\ref{tab:demo_users}
liste les identifiants utilisés en environnement de développement.

\begin{table}[H]
\caption{Comptes de démonstration (mot de passe : \texttt{password})}
\label{tab:demo_users}
\centering
\small
\begin{tabular}{|l|l|p{5.5cm}|}
\hline
\textbf{Email} & \textbf{Profil} & \textbf{Accès principal} \\
\hline
\texttt{admin@civco.ma} & Administrateur & Configuration, tous modules \\
\hline
\texttt{ingenieur@civco.ma} & Chef de projet & Projets, tâches, équipe \\
\hline
\texttt{compta@civco.ma} & Comptable & Factures, paiements \\
\hline
\texttt{chantier@civco.ma} & Chef de chantier & Projets, tâches terrain \\
\hline
\texttt{tech@civco.ma} & Technicien & Tâches assignées, consultation \\
\hline
\end{tabular}
\end{table}

\medskip

Ces comptes sont créés par la commande \texttt{php artisan migrate --seed} et
permettent de valider le contrôle d'accès par rôles décrit au Chapitre~3.

% =====================================================================
\section{Référentiel des endpoints API}

Le Tableau~\ref{tab:api_endpoints} présente les principaux endpoints exposés
sous le préfixe \texttt{/api/v1/}. Toutes les routes (sauf \texttt{login}) sont
protégées par le middleware \texttt{auth:sanctum} et soumises à un contrôle de
permissions.

\begin{table}[H]
\caption{Principaux endpoints de l'API REST}
\label{tab:api_endpoints}
\centering
\footnotesize
\begin{tabular}{|l|l|p{5.5cm}|}
\hline
\textbf{Méthode} & \textbf{Endpoint} & \textbf{Description} \\
\hline
POST & \texttt{/login} & Authentification utilisateur \\
\hline
GET & \texttt{/me} & Profil et permissions de l'utilisateur courant \\
\hline
GET/POST & \texttt{/projects} & Liste et création de projets \\
\hline
GET/PUT & \texttt{/projects/\{id\}} & Consultation et modification d'un projet \\
\hline
GET/POST & \texttt{/clients} & Liste et création de clients \\
\hline
GET/POST & \texttt{/quotes} & Liste et création de devis \\
\hline
POST & \texttt{/quotes/\{id\}/convert} & Conversion devis $\rightarrow$ facture \\
\hline
GET/POST & \texttt{/invoices} & Liste et création de factures \\
\hline
GET/POST & \texttt{/tasks} & Liste et création de tâches \\
\hline
GET/POST & \texttt{/contracts} & Liste et création de contrats \\
\hline
GET & \texttt{/dashboard} & Indicateurs du tableau de bord \\
\hline
GET & \texttt{/search} & Recherche globale (clients, projets, documents) \\
\hline
GET & \texttt{/activity-logs} & Journal des actions (administrateur) \\
\hline
\end{tabular}
\end{table}

% =====================================================================
\section{Schéma des entités principales}

Le Tableau~\ref{tab:entity_schema} décrit les attributs essentiels des entités
métier les plus consultées dans l'application.

\begin{table}[H]
\caption{Attributs principaux des entités métier}
\label{tab:entity_schema}
\centering
\small
\begin{tabular}{|l|p{10.5cm}|}
\hline
\textbf{Entité} & \textbf{Attributs clés} \\
\hline
\texttt{Project} & \texttt{reference}, \texttt{title}, \texttt{status}, \texttt{budget},
\texttt{progress\_percent}, \texttt{start\_date}, \texttt{end\_date},
\texttt{latitude}, \texttt{longitude}, \texttt{client\_id} \\
\hline
\texttt{Client} & \texttt{name}, \texttt{email}, \texttt{phone}, \texttt{status},
\texttt{is\_official}, \texttt{archived\_at} \\
\hline
\texttt{Quote} & \texttt{reference}, \texttt{status}, \texttt{total\_ht},
\texttt{total\_tax}, \texttt{total\_ttc}, \texttt{project\_id}, \texttt{client\_id} \\
\hline
\texttt{Invoice} & \texttt{reference}, \texttt{status}, \texttt{issued\_at},
\texttt{due\_date}, \texttt{amount\_paid}, \texttt{balance\_due}, \texttt{quote\_id} \\
\hline
\texttt{Task} & \texttt{title}, \texttt{status}, \texttt{priority},
\texttt{due\_date}, \texttt{project\_phase\_id}, \texttt{assigned\_to} \\
\hline
\texttt{Contract} & \texttt{reference}, \texttt{status}, \texttt{signed\_at},
\texttt{project\_id}, \texttt{template\_id} \\
\hline
\end{tabular}
\end{table}

\medskip

Les statuts des entités sont implémentés sous forme d'énumérations PHP
(\texttt{ProjectStatus}, \texttt{QuoteStatus}, \texttt{InvoiceStatus},
\texttt{TaskStatus}) garantissant la validité des transitions d'état côté
serveur.

% =====================================================================
\section{Guide de démarrage rapide}

Pour reproduire l'environnement de développement décrit au Chapitre~4, les commandes
suivantes permettent d'initialiser l'application à partir du dépôt source :

\begin{lstlisting}[language={}, caption={Commandes d'initialisation de l'application}]
# Backend
cd backend
composer install
cp .env.example .env
php artisan key:generate
php artisan migrate --seed
php artisan serve

# Frontend (dans un second terminal)
cd frontend
npm install
npm run dev
\end{lstlisting}

\medskip

L'application est alors accessible à l'adresse \texttt{http://127.0.0.1:5173} (frontend)
avec le proxy Vite redirigeant les appels API vers \texttt{http://127.0.0.1:8000}.
Les comptes de démonstration listés en section~3 permettent de tester les différents
profils métier.

\newpage
